\pagestyle{index}
\phantomsection\addcontentsline{toc}{section}{Index des auteurs}
\hypertarget{indexa}{}
\printindex[a]

\clearpage
\pagestyle{plain}
\cleardoublepage
\pagestyle{index}

\phantomsection\addcontentsline{toc}{section}{Index des textes selon leur genre}
\hypertarget{indexg}{}
\printindex[g]

\clearpage
\pagestyle{plain}
\cleardoublepage
\pagestyle{index}

\begin{spacing}{0.92}
\phantomsection\addcontentsline{toc}{section}{Index des questions selon les axes du programme puis selon leur type}
\hypertarget{indexw}{}
\printindex[w]
\end{spacing}

\clearpage
\pagestyle{plain}
\cleardoublepage
\pagestyle{index}

\begin{spacing}{0.92}
\phantomsection\addcontentsline{toc}{section}{Index des questions selon leur type puis selon les axes du programme}
\hypertarget{indexv}{}
\printindex[v]
\end{spacing}

\clearpage
\pagestyle{plain}
\cleardoublepage
\pagestyle{index}

\begin{spacing}{0.7}
\phantomsection\addcontentsline{toc}{section}{Index des questions selon leur type}
\hypertarget{indexq}{}
\printindex[q]
\end{spacing}

%\clearpage
%\pagestyle{plain}
%\cleardoublepage
%\pagestyle{index}
%
%\begin{spacing}{0.7}
%\phantomsection\addcontentsline{toc}{section}{Index questions selon les axes du programme}
%\hypertarget{indexp}{}
%\printindex[p]
%\end{spacing}

\clearpage
\pagestyle{plain}
\cleardoublepage
\pagestyle{index}

\phantomsection\addcontentsline{toc}{section}{Index thématique}
%\label{indext}
\printindex[t]

\stepcounter{chapter}

\section*{Se repérer dans l'index thématique}
\phantomsection\addcontentsline{toc}{subsection}{Se repérer dans l'index thématique}\hypertarget{repere-thematique}{}\label{repere-thematique}\stepcounter{section}

Parce qu'il peut être difficile, au vu du nombre de termes indexés, de se repérer pour trouver ce que l'on cherche, et parce qu'il est toujours difficile de savoir quoi indexer, on propose ci-dessous une liste des grandes thématiques par chapitre et des termes que l'on a indexés sous ces thématiques. Cela met aussi en valeur les entrées sur lesquelles les thématiques se croisent.

\subsection{Éducation, transmission et émancipation}
\begin{itemize}\setlength{\itemsep}{6pt}

\item \textbf{Sur la difficulté à transmettre~:} apprentissage, connaissances, éducation, ennui, innocence, parole.

\item \textbf{Sur la transmission familiale~:} ambition, autorité, culture, enfant, \textit{la sous-entrée \emph{et éducation} de l'entrée} famille, identification, transmission.

\item \textbf{Sur l'apprentissage de la liberté~:} citoyen, éducation, ennui, facultés, formation, innocence, liberté.

\item \textbf{Sur le rapport à la culture~:} ambition, culture, lecture, loisirs.

\item \textbf{Sur la question des apprentissages~:} apprentissage, châtiments, curiosité, éducation \textit{et ses sous-entrées}, enfant, explications, ignorance, instruction, leçon, parole, préjugés, savoir, \textit{la sous-entrée \emph{et éducation} de l'entrée} société.

\end{itemize}

%%%%%%%%%%%%%%%Chapitre 2%%%%%%%%%%%%%%%

\subsection{Expressions de la sensibilité}

\begin{itemize}\setlength{\itemsep}{6pt}

\item \textbf{Autour des émotions~:} amour, deuil, émotions, larmes, passion, pitié sensibilité, sentiments, souffrance, solitude, tristesse.

\item \textbf{Sur l'art (et la création artistique)~:} art, beauté, contemplation, création, désir, divertissement, expérience esthétique, inspiration, intellection, passion, sensation, sensibilité, solitude, sublime.

\item \textbf{Au sujet du rapport de l'humain à la nature dans le romantisme~:} nature (environnement), nuit, paysage, sensibilité.

\item \textbf{Autour du son~:} bruit, musique, silence, solitude, son.

\item \textbf{Sur la sensibilité aux autres ou aux animaux~:} amitié, conscience, morale, pitié, sensibilité, sentiments, vertu.

\end{itemize}

%%%%%%%%%%%%%%%Chapitre 3%%%%%%%%%%%%%%%

\subsection{Métamorphoses du moi}

\begin{itemize}\setlength{\itemsep}{6pt}

\item \textbf{Sur le rapport à soi~:} biographie, changement, conscience, corps, intime, introspection \textit{et sa sous-entrée, « se trouver »}, langage, mensonge, mémoire, moi \textit{et ses sous-entrées}, narcissisme, personnage, personnalité, sincérité, soi \textit{et ses sous-entrées}, solitude, sujet, tristesse.

\item \textbf{Sur l'identité~:} apparences, caractère, changement, corps, identité \textit{et ses sous-entrées}, intériorité, masque, mémoire, passé, personnage, personnalité, rêve (problème du rêve et du sommeil (Descartes)), singularité, substance, sujet, vie et mort (âge de la vie).

\item \textbf{Sur le rapport à l'autre~:} autrui \textit{et ses sous-entrées}, famille, personnage, reconnaissance, rôle (jouer un rôle), société (organisation sociale), théâtre.

\item \textbf{Sur les âges ou étapes de la vie~:} adulte, adolescence,  enfant, enfance, existence, famille, maladie, naissance, parent, regrets, souvenir, temps, vie (âge de la vie), vieillesse.

\item \textbf{Sur la morale~:} bonheur, désir, divertissement, imagination, morale \textit{et sa sous-entrée conscience morale}, passion.

\item \textbf{Sur la conscience et l'inconscient~:} conscience, inconscient, pulsions, psychanalyse, rêve.

\item \textbf{Sur le récit de soi~:} communication, écriture, enfance, postérité, relation.

\end{itemize}

%%%%%%%%%%%%%%%Chapitre 4%%%%%%%%%%%%%%%
\stepcounter{section}
\subsection{Création, continuités et ruptures}

\begin{itemize}\setlength{\itemsep}{6pt}

\item \textbf{Sur le progrès scientifique, et les éventuelles discontinuités~:} électricité, lois de la nature, sciences, scientifique, société, technique \textit{et sciences}, technologie.

\item \textbf{Sur les ruptures artistiques, la nouveauté et les avant-gardes}~: ancien, architecture, art et \textit{sa sous-entrée style artistique}, avant-garde, cubisme, culture, idéal, morale, nouveauté, passion, progrès, \textit{la sous-entrée \emph{religion} de l'entrée} art.

\item \textbf{Sur la création littéraire~:} création, critique (littéraire), écrivain, génération, génie, influence, lecture, nouveauté, penseur, tradition.

\end{itemize}

%%%%%%%%%%%%%%%Chapitre 5%%%%%%%%%%%%%%%

\subsection{Histoire et violence}

\begin{itemize}\setlength{\itemsep}{6pt}

\item \textbf{Sur la Shoah~:} camp de concentration, centre de mise à mort, déportation, exécution, guerre (Seconde Guerre mondiale), Résistants, Shoah, SS, témoignage, traumatisme, violence.

\item \textbf{Sur la guerre en général~:} armée, armes, bataille, blessure, brutalisation, classes sociales, combat, esthétique de la violence, exécution, folie, guerre \textit{et ses sous-entrées}, marche (militaire), mort, paix, patriotisme, poilus, religion, stratégie, torture, traumatisme, valeurs, victime, vie (et mort), violence.

\item \textbf{Au sujet de la capacité à témoigner du passé ~:} absurde, enfant, incompréhension, mémoire, oubli, souvenir, témoignage \textit{et sa sous-entrée}, traumatisme, violence.

\item \textbf{Sur le rapport entre État et violence~:} démocratie, désobéissance, État, idéologie, individu, injustice, justice, patriotisme, pensée \textit{et ses sous-entrées}, politique, pouvoir, religion, sacrifice, violence.

\item \textbf{Sur l'industrialisation et la vie ouvrière~:} aliénation, artificiel, condition ouvrière, corps, dialectique, dignité humaine, déshumanisation, machines, usine, violence.

\item \textbf{Sur l'histoire~:} histoire, sacrifice, temps, tradition, transmission.

\item \textbf{Sur la (dé)colonisation~:} anthropocentrisme, civilisation, colonisation, commerce, culture, voyage.

\item \textbf{Sur la paix~:} confiance, désarmement, droit \textit{et sa sous-entrée} droit international, \textit{la sous-entrée} relations inter-étatiques \textit{de l'entrée} État, paix, \textit{et sa sous-entrée} pacifisme.

\end{itemize}

%%%%%%%%%%%%%%%Chapitre 6%%%%%%%%%%%%%%%

\subsection{L'humain et ses limites}

\begin{itemize}\setlength{\itemsep}{6pt}

\item \textbf{Sur la transformation technique de l'humain (transhumanisme)~:} anomalie, artificiel, biologie, biogénétique, déshumanisation, eugénisme, monstre, nature (nature humaine), technique \textit{et ses sous-entrées}, transhumanisme.

\item \textbf{Sur la finitude humaine~:} finitude, \textit{la sous-entrée \emph{âge de la vie} de l'entrée} vie, vieillissement.

\item \textbf{Sur l'écologie~:} écologie, nature (environnement), principe de précaution, technique (maîtrise technique de la nature).

\item \textbf{Sur la déshumanisation ou le posthumanisme (foucaldien)~:} déshumanisation, eugénisme, individu, monstre, nature (nature humaine), société \textit{et ses sous-entrées)}, valeurs, violence \textit{et ses sous-entrées}.

\item \textbf{Sur l'influence de la société sur l'humanité~:}
dépersonnalisation, famille, individu, individualisme, société \textit{et ses sous-entrées}, technique \textit{et ses sous-entrées}.

\end{itemize}

\vfill{~}

\clearpage
\pagestyle{plain}
\cleardoublepage
\pagestyle{index}

\begin{spacing}{1}
\phantomsection\addcontentsline{toc}{section}{Index des sujets par lieux et années}
\hypertarget{indexn}{}
\indexprologue{\textit{Sauf mention du contraire, les épreuves mentionnées ci-dessous étaient les épreuves de la session « normale » ou « obligatoire » du baccalauréat. Pour les sujets des épreuves qui ont eu lieu lors de la session de septembre (session de remplacement) ou destinés aux candidats libres, quand de telles épreuves ont eu lieu, on l'a précisé.}}%
\printindex[n]
\end{spacing}